\documentclass{article}

% Language setting
\usepackage[english]{babel}

% Set page size and margins
\usepackage[letterpaper,top=2cm,bottom=2cm,left=3cm,right=3cm,marginparwidth=1.75cm]{geometry}

% Useful packages
\usepackage{amsmath}
\usepackage{amssymb}
\usepackage{graphicx}
\usepackage[colorlinks=true, allcolors=black]{hyperref}
\usepackage{titlesec}
\usepackage{xparse}

\NewDocumentCommand{\qcq}{m o o m}{%
	\ensuremath{\forall #1 \IfValueT{#2}{, #2} \IfValueT{#3}{, #3}  \in  \mathbb{#4}}%
}
\NewDocumentCommand{\ext}{m o o m}{%
	\ensuremath{\exists #1 \IfValueT{#2}{, #2} \IfValueT{#3}{, #3}  \in  \mathbb{#4}}%
}

\title{\textbf{Exercices : Calcul Intégral}}
\author{\textbf{Yassin Akkab}}

\newcommand{\R}{\mathbb{R}}
\newcommand{\N}{\mathbb{N}}
\newcommand{\Z}{\mathbb{Z}}
\newcommand{\C}{\mathbb{C}}
\newcommand{\Q}{\mathbb{Q}}
\newcommand{\D}{\mathbb{D}}

\begin{document}
	
	\maketitle
	
	\subsection*{Exercice 1 (Calcul de primitives et intégrales simples)}
	Calculez les intégrales suivantes :
	\[ I_1 = \int_0^1 \frac{x}{x^2 + 1} \, dx \]
	\[ I_2 = \int_0^{\pi/2} \sin(x)\cos^3(x) \, dx \]
	\[ I_3 = \int_1^e \frac{\ln(x)}{x} \, dx \]
	
	\subsection*{Exercice 2 (Intégration par parties)}
	Calculez les intégrales suivantes à l'aide d'une intégration par parties :
	\[ J_1 = \int_1^e x \ln(x) \, dx \]
	\[ J_2 = \int_0^\pi x \sin(x) \, dx \]
	\[ J_3 = \int_0^1 \arctan(x) \, dx \]
	
	\subsection*{Exercice 3 (Sommes de Riemann)}
	Soit la suite $(u_n)_{n \in \mathbb{N}^*}$ définie par :
	\[ u_n = \sum_{k=1}^n \frac{n}{n^2 + k^2} \]
	\begin{enumerate}
		\item Montrer que pour tout $n \ge 1$ :
		\[ u_n = \frac{1}{n} \sum_{k=1}^n \frac{1}{1 + \left(\frac{k}{n}\right)^2} \]
		\item En utilisant le théorème des sommes de Riemann pour une fonction $f$ continue à déterminer sur $[0, 1]$, calculer la limite de la suite $(u_n)$.
	\end{enumerate}
	
	\subsection*{Exercice 4 (Étude d'une fonction définie par une intégrale)}
	Soit $F$ la fonction définie sur $[0, +\infty[$ par :
	\[ F(x) = \int_0^x e^{-t^2} \, dt \]
	\begin{enumerate}
		\item Montrer que $F$ est dérivable sur $[0, +\infty[$ et calculer sa dérivée $F'$.
		\item Étudier le sens de variation de $F$.
		\item Montrer que pour tout $t \ge 1$, $e^{-t^2} \le e^{-t}$.
		\item En déduire que pour tout $x \ge 1$ :
		\[ F(x) \le F(1) + e^{-1} - e^{-x} \]
		\item En déduire que la fonction $F$ est majorée sur $[0, +\infty[$, et qu'elle admet une limite finie en $+\infty$.
	\end{enumerate}
	
\end{document}
